
\section{系统应用示例}

\subsection{应用场景与使用角色}

本系统在实际应用中主要对应三类用户角色。第一类是专利计量研究者，关注分析流程的可复现性、变量定义的可解释性以及统计结论的可验证性。第二类是技术情报分析师，关注快速形成多视角研判并输出可执行洞察。第三类是企业知识产权管理者，关注高价值技术布局识别、竞争风险监测与专利资产配置。

与传统手工流程相比，系统将"问题理解$\rightarrow$方案设计$\rightarrow$代码执行$\rightarrow$结果解释"串联为自动化闭环，尤其在规划阶段引入数据证据，降低了人工试错成本。对上述三类角色而言，系统的核心价值不是替代领域判断，而是提供稳定、可追踪、可复核的分析底座。

为展示系统在不同分析场景下的应用方式与输出形态，本章选取三个研究问题作为应用案例，分别对应因果驱动分析、技术融合趋势识别和国家竞争态势评估三类典型专利计量任务。三者共同覆盖“机制解释”“结构趋势识别”“竞争比较评估”三种常见研究目标，因此可作为系统能力展示的说明性案例集。三个案例均采用完整系统（C\_iterative 模式），所有数据来自实际实验产出文件。

需要说明的是，本章定位为\textbf{流程展示与中间产物说明}，重点回答“系统如何把研究问题转化为分步任务、代码和结果解释”；其作用是补充展示第五章定量实验背后的执行链路，而非提供独立于第五章之外的新一轮效果验证。因此，本章中的案例结果应被理解为完整系统在三个代表性问题上的样例输出，系统优劣的统一判断仍以第五章的对比实验为准。

\subsection{案例一：技术影响力驱动因素分析}

本案例选择 Q1"分析数据安全领域的技术影响力驱动因素"，展示从输入目标到输出报告的两轮完整分析过程\footnote{本案例数据来源为系统实验产出文件 \texttt{experiment\_q1\_C\_iterative.json}。}。

\subsubsection{输入与数据感知}

用户输入研究问题后，系统解析目标类型和数据边界。研究目标为"分析数据安全领域的技术影响力驱动因素"，目标类型为因果/驱动因素识别，数据集为 \texttt{data/new\_data.XLSX}，包含 10,000 条数据安全领域专利、22 个字段。

Strategist 调用 DataPreview 与 GraphPreview，生成 3 条被选中用于蓝图规划的洞察：

\textbf{洞察 I1}（statistical）：近年授权专利被引量骤降——反映技术迭代加速，需在分析中控制专利年龄，否则会混淆真实影响力。

\textbf{洞察 I2}（graph\_structural）：H04L9/40 为核心枢纽——可限定分析范围至该 IPC 主导的子网络，提升假设检验的针对性。

\textbf{洞察 I3}（statistical）：法律状态复杂性提升引用活跃度——提示需控制法律状态以隔离商业化干扰。

\subsubsection{第一轮分析：探索性验证}

基于洞察和双知识图谱，Strategist 生成第一轮 DAG 蓝图，包含 2 个分析任务：

\textbf{任务 1：变量计算（variable\_calculation）}

% 请确保导言区包含：\usepackage{booktabs, array}

\begin{table}[htbp]
\centering
\caption{案例一第一轮任务1配置}
\label{tab:case1-r1-t1}
\begin{tabular}{>{\centering\arraybackslash}p{3.2cm} >{\raggedright\arraybackslash}p{10.4cm}}
\toprule % 顶线
项目 & 内容 \\
\midrule % 栏目线
研究问题 & 如何计算每项专利的技术多样性并标准化其影响力以排除时间偏差？ \\
使用列   & IPC分类号、被引用专利数量、授权日 \\
参数     & reference\_year=2026 \\
\bottomrule % 底线
\end{tabular}
\end{table}

执行结果：有效专利 9,879 条，技术多样性均值为 3.89（中位数 3.0，标准差 2.88，范围 1-78），呈明显右偏分布。以 2026 年 1 月为基准时点，平均专利年龄为 2.7 年，标准化被引量均值为 0.57，有效消除了专利年龄带来的时间累积效应，与洞察 I1 的判断一致。

\textbf{任务 2：假设检验（hypothesis\_test）}

% 请确保在 LaTeX 文档的导言区添加：
% \usepackage{booktabs, array}

\begin{table}[htbp]
\centering
\caption{案例一第一轮任务2配置}
\label{tab:case1-r1-t2}
\begin{tabular}{>{\centering\arraybackslash}p{3.2cm} >{\raggedright\arraybackslash}p{10.4cm}}
\toprule % 顶线（粗线）
项目 & 内容 \\
\midrule % 栏目线（细线）
研究问题 & 技术跨界度对技术影响力是否有负向影响？ \\ \addlinespace[0.5ex] % 增加一点行间距，防止文字拥挤
使用列   & IPC分类号、被引用专利数量、授权日、法律状态/事件、IPC主分类号 \\ \addlinespace[0.5ex]
控制变量 & patent\_age, legal\_status\_dummies \\
\bottomrule % 底线（粗线）
\end{tabular}
\end{table}

执行结果：假设未获支持。全样本回归中技术多样性系数不显著（$\beta$=0.006, p=0.909, R²=0.066, n=10,000）；在 H04L9/40 子群中系数方向为负但同样不显著（$\beta$=-0.157, p=0.147, R²=0.039, n=816）。该"零结果"否定了"技术融合必然损害影响力"的假设，为第二轮提供了明确的探索方向。

\subsubsection{第二轮分析：差距修补与深化验证}

系统读取第一轮结果后，自动进行差距分析：第一轮排除了技术跨界度作为显著驱动因素的假设，但洞察中提及的法律状态因果关系、H04L9/40 共现强度影响、高连接度发明人等方向未被检验。系统据此在第二轮引入控制变量（授权年份、IPC 主分类号前 4 位）、交互项（法律状态×IPC 核心度）及非线性建模（负二项回归）。

第二轮 DAG 蓝图包含 2 个深化任务：

\textbf{任务 R2\_1：IPC 共现网络特征工程 + 负二项回归}

% 请确保导言区已添加：\usepackage{booktabs, array}

\begin{table}[htbp]
\centering
\caption{案例一第二轮任务R2\_1配置}
\label{tab:case1-r2-t1}
\begin{tabular}{>{\centering\arraybackslash}p{3.2cm} >{\raggedright\arraybackslash}p{10.4cm}}
\toprule % 顶线
项目 & 内容 \\
\midrule % 栏目线
研究问题 & IPC 共现网络中 H04L9/40 的共现强度是否显著正向预测专利被引量？ \\ \addlinespace[0.5ex]
方法     & 负二项回归（替代 OLS，适应计数数据） \\ \addlinespace[0.5ex]
控制变量 & 授权年份固定效应, IPC 主分类前缀 \\
\bottomrule % 底线
\end{tabular}
\end{table}

该任务将 GraphPreview 的 IPC 共现网络信息转化为回归自变量，实现"图特征$\rightarrow$统计建模"路径。

\textbf{任务 R2\_2：发明人网络分析 + 文本挖掘 + 中介检验}

% 请确保在导言区添加：\usepackage{booktabs, array}

\begin{table}[htbp]
\centering
\caption{案例一第二轮任务R2\_2配置}
\label{tab:case1-r2-t2}
\begin{tabular}{>{\centering\arraybackslash}p{3.2cm} >{\raggedright\arraybackslash}p{10.4cm}}
\toprule % 顶线
项目 & 内容 \\
\midrule % 栏目线
研究问题 & 高连接度发明人所参与的专利是否更倾向于包含特定技术主题，且该主题是否中介其对被引量的正向影响？ \\ \addlinespace[0.5ex]
方法     & BERT 主题建模 + 负二项回归 + 中介效应检验 \\ \addlinespace[0.5ex]
参数     & high\_degree\_threshold=30 \\
\bottomrule % 底线
\end{tabular}
\end{table}

该任务同时利用发明人合作网络和摘要内容，实现"网络位置$\rightarrow$技术主题$\rightarrow$影响力"的跨层级机制检验。

两轮合计产出 4 个分析任务，全部执行成功。分析层次从第一轮的"效应是否存在"提升至第二轮的"效应的机制、路径和边界条件"。

\subsection{案例二：技术融合趋势识别}

本案例选择 Q2"识别数据安全领域的技术融合趋势"，重点展示系统对图结构信息和文本内容的综合利用能力\footnote{本案例数据来源为系统实验产出文件 \texttt{experiment\_q2\_C\_iterative.json}。}。

\subsubsection{数据感知与洞察选择}

系统为本问题选中 3 条核心洞察：

\textbf{洞察 I1}（graph\_structural）：H04L9/40 在 IPC 共现网络中度数高达 904，是技术融合网络的核心枢纽，可作为融合分析的锚点。

\textbf{洞察 I2}（statistical）：质押状态专利平均被引 84.08 次，远高于普通授权专利的 11.19 次，提示高融合度专利具有更高商业价值。

\textbf{洞察 I3}（content\_driven）：摘要中频繁出现"安全组件嵌入"、"SoC 测试密钥"等硬件级安全术语，结合 IPC 分布跨越 A、F、G 等部类，提示硬件级安全机制成为跨领域融合新焦点。

\subsubsection{第一轮分析：融合结构探索}

第一轮 DAG 蓝图包含 3 个任务，分别从网络结构、法律状态关联和时间趋势三个角度探索融合特征：

\textbf{任务 1：H04L9/40 桥接 IPC 提取（variable\_calculation）}。从 IPC 共现网络中提取 H04L9/40 子图，识别直接连接的 IPC 编码并按部类聚合，计算跨部类边占比。结果显示 H04L9/40 子图中跨部类边占比达 78.3\%，Top 3 桥接分类为 G16H50/70（医疗 AI）、A61B5/00（医疗设备）和 B61L27/00（铁路控制），表明安全协议已深度渗透至医疗和交通领域。

\textbf{任务 2：法律状态与 IPC 多样性关联检验（hypothesis\_test）}。按法律状态分组计算 IPC Shannon 熵和跨部类引用比例，进行独立样本 t 检验。结果显示质押专利 IPC Shannon 熵为 2.91，显著高于普通授权专利的 2.43（p<0.01，效应量 d=0.42），跨部类引用比例高出 31\%，提示技术融合度与商业价值之间存在显著相关性。

\textbf{任务 3：硬件级安全专利时间趋势（variable\_calculation）}。使用 BERT 模型对摘要进行安全层级分类（硬件/固件/软件），按年度聚合。结果显示，硬件安全专利占比由 2023 年的 18.2\% 提升至截至 2026 年 1 月样本的 34.7\%。由于 2026 年仅覆盖截至 1 月的最新观测点，本文不将其解释为完整年度增速，而仅比较同口径下的阶段性升幅；按该口径，硬件层的上升幅度仍高于固件层（+9.3\%）和软件层（-5.2\%），说明硬件级安全议题持续升温。

\subsubsection{第二轮分析：融合机制深化}

系统识别第一轮的差距：未对 H04L9/40 共现专利进行语义层面的主题建模，未量化硬件安全关键词在各 IPC 部类的渗透强度差异，法律状态对引用的影响未控制 IPC 多样性作为混杂变量。第二轮生成 3 个深化任务：

\textbf{任务 R2\_1：层次回归分析}。在控制 IPC 分类多样性和技术子域固定效应后，质押专利仍显示额外 62.3 次引用增量（p<0.001），权利转移状态额外增加 18.7 次（p=0.003），证实法律状态复杂性独立于技术广度对引用深度的正向影响。

\textbf{任务 R2\_2：跨领域硬件安全术语提取}。使用预训练中文 BERT 提取摘要中的硬件安全术语，按 IPC 部类分组计算年度关键词密度。结果显示，A/F/G/H 部类的硬件安全关键词密度持续上升，截至 2026 年 1 月的最新观测点仍明显高于其他部类。按同一阶段口径估算，A/F/G/H 部类关键词密度增幅为 37.2\%，高于其他部类的 8.1\%；其中 A 部（医疗）和 F 部（机械）增幅最突出，分别为 52.4\% 和 48.7\%。

\textbf{任务 R2\_3：IPC-摘要联合主题建模}。对含 H04L9/40 的专利进行 LDA 主题建模（K=8），结合共现 IPC 标注语义。识别出 5 个显著融合簇：5G+医疗（28.3\%）、车联网认证（22.1\%）、工业物联网安全（19.7\%）、区块链身份（15.4\%）和消费电子可信执行环境（14.5\%）。

两轮合计产出 6 个分析任务，全部执行成功。系统从宏观融合结构（跨部类边占比）逐步深入到微观融合簇识别（5 个具体场景），体现了两轮迭代从"现象发现"到"机制解释"的分析深化路径。

\subsection{案例三：国家竞争态势评估}

本案例选择 Q3"评估不同国家/地区在数据安全领域的竞争态势"，重点展示系统在比较分析和多维指标构建场景下的能力\footnote{本案例数据来源为系统实验产出文件 \texttt{experiment\_q3\_C\_iterative.json}。}。

\subsubsection{数据感知与洞察选择}

系统为本问题选中 3 条核心洞察：

\textbf{洞察 I1}（statistical）：质押专利被引量异常显著（84.08 vs 11.19），可作为"高价值专利"的量化标识，避免仅依赖数量指标评估竞争力。

\textbf{洞察 I2}（graph\_structural）：美国在优先权合作网络中占据核心位置（度数=14），反映其在全球技术部署中的枢纽地位。

\textbf{洞察 I3}（graph\_structural）：H04L9/40 作为核心技术锚点，可聚焦分析各国在该协议上的布局差异，避免泛化的"数据安全"讨论。

\subsubsection{第一轮分析：竞争格局探索}

第一轮 DAG 蓝图包含 3 个任务：

\textbf{任务 1：质押专利与网络中心性关联检验（hypothesis\_test）}。按法律状态分组计算平均入度（被引次数），进行双样本 t 检验。质押专利平均被引 84.08 次，显著高于普通授权专利的 11.19 次（Cohen's d=0.67，p<0.00001），支持将质押状态视为高价值专利的候选标识。

\textbf{任务 2：国家网络中心性与产出质量回归分析（regression）}。构建优先权国家合作网络，计算各国度中心性，与专利族规模和引用影响力进行回归分析。国家度中心性对族规模和引用影响力呈中等正向效应（$\beta$=0.42，p=0.001，R²=0.31），美国（度数=14）在技术产出数量和全球控制力上均显著领先。

\textbf{任务 3：H04L9/40 核心协议的国家分布分析（data\_summary）}。按优先权国家对 H04L9/40 主分类专利进行分组统计。美国占 38\%（平均 4.2 次引用），侧重"动态证书生成"；中国占 25\%（平均 2.1 次引用），侧重"端到端加密"；韩国占 12\%，侧重"5G 用户面安全"。各国在核心协议上呈现差异化技术路线。

\subsubsection{第二轮分析：竞争机制深化}

系统识别第一轮的差距：未系统整合国家、技术、价值三维关系，未对 H04L9/40 内部子领域进行国家级主题聚类，质押专利的高引用行为未控制专利年龄的混杂影响。第二轮生成 2 个深化任务：

\textbf{任务 R2\_1：国家-技术子类聚类与优势分析}。对含 H04L9/40 的专利提取剩余 IPC 编码和摘要文本，进行主题聚类，计算各簇的国家分布 Shannon 熵和主导国家占比，识别各国技术优势领域。

\textbf{任务 R2\_2：高价值专利国家归属与年龄控制引用分析}。提取质押状态样本标注优先权国家，计算专利年龄（2026 - 授权年），使用协方差分析（ANCOVA）以被引量为因变量、法律状态为因素、专利年龄为协变量，检验质押组的引用优势是否独立于年龄效应，并报告质押专利国家 Top 5 分布。

两轮合计产出 5 个分析任务，全部执行成功。系统从第一轮的宏观竞争格局描述（国家数量对比、中心性排名）深化至第二轮的微观竞争机制解释（子领域优势识别、年龄控制后的独立效应验证）。

\subsection{三案例综合评价}

\subsubsection{执行统计对比}

% 请确保在导言区添加：\usepackage{booktabs, array}

\begin{table}[htbp]
\centering
\caption{三案例执行统计对比}
\label{tab:case-comparison}
\resizebox{\textwidth}{!}{
\begin{tabular}{>{\centering\arraybackslash}p{2.3cm} >{\raggedright\arraybackslash}p{3.8cm} >{\raggedright\arraybackslash}p{3.8cm} >{\raggedright\arraybackslash}p{3.8cm}}
\toprule % 顶线
维度 & Q1 影响力驱动因素 & Q2 技术融合趋势 & Q3 国家竞争态势 \\
\midrule % 栏目线
第一轮任务数 & 2 & 3 & 3 \\
第二轮任务数 & 2 & 3 & 2 \\
任务总数     & 4 & 6 & 5 \\
执行成功率   & 100\% & 100\% & 100\% \\ \addlinespace[0.5ex]
分析类型覆盖 & 回归、中介检验 & 图分析、主题建模、层次回归 & 网络分析、ANCOVA、聚类 \\
\bottomrule % 底线
\end{tabular}
}
\end{table}

三个案例在任务结构上存在显著差异：Q1 偏因果识别，任务链为"变量构建$\rightarrow$假设检验$\rightarrow$机制验证"；Q2 偏结构与趋势，任务链为"网络提取$\rightarrow$关联检验$\rightarrow$主题聚类"；Q3 偏比较与评估，任务链为"指标构建$\rightarrow$国家对比$\rightarrow$优势归因"。这一差异与系统"共框架、异任务"的设计目标相符——同一架构能够根据研究问题自动生成差异化的分析路径。

\subsubsection{两轮迭代效果对比}

% 请确保导言区已添加：\usepackage{booktabs, array}

\begin{table}[htbp]
\centering
\caption{两轮迭代效果对比}
\label{tab:iteration-comparison}
\resizebox{\textwidth}{!}{
\begin{tabular}{>{\centering\arraybackslash}p{2.3cm} >{\raggedright\arraybackslash}p{5.7cm} >{\raggedright\arraybackslash}p{5.7cm}}
\toprule % 顶线
维度 & 第一轮共性特征 & 第二轮共性特征 \\
\midrule % 栏目线
分析层次     & 效应存在性检验 & 机制、路径与边界条件 \\ \addlinespace[0.5ex]
典型方法     & OLS 回归、t 检验、描述统计 & 负二项回归、层次回归、ANCOVA、主题建模 \\ \addlinespace[0.5ex]
数据源       & 结构化字段 & 结构化字段 + 图特征 + 文本内容 \\ \addlinespace[0.5ex]
控制变量     & 少量（1-2 个） & 多层次（年龄、固定效应、交互项） \\ \addlinespace[0.5ex]
典型深化模式 & — & 非线性替代、混杂控制、跨层级机制检验 \\
\bottomrule % 底线
\end{tabular}
}
\end{table}

三个案例均展现了从第一轮到第二轮的一致性深化模式：第一轮建立基础变量关系并发现初步信号（包括零结果），第二轮基于差距分析引入更严格的统计控制、更丰富的数据源和更深层的机制检验。该模式表明两轮迭代机制在这三类分析场景下具有较好的适用性。

\subsubsection{应用效果总结}

基于三个案例，可总结三方面应用效果：

第一，方案可解释性。每个任务均可追溯至洞察条目和图谱路径。例如 Q1 任务 1 的"标准化被引量消除时间偏差"源于洞察 I1，Q1 任务 2 的"控制法律状态"源于洞察 I3，Q2 任务 1 的"H04L9/40 子图提取"源于洞察 I1，Q3 任务 2 的"国家度中心性"源于洞察 I2。这种"洞察$\rightarrow$任务$\rightarrow$结论"的证据链避免了黑箱式自动分析。

第二，执行可复核性。三个案例共 15 个任务全部执行成功，蓝图 JSON、技术规格、Python 代码和执行结果均有落盘文件，可复跑、可对比、可审计。

第三，迭代驱动深化。三个案例均展现了"第一轮发现$\rightarrow$差距识别$\rightarrow$第二轮深化"的自适应模式。Q1 从"零结果"转向替代机制探索，Q2 从宏观融合结构深入到 5 个具体融合簇，Q3 从国家数量对比深化至年龄控制后的独立效应验证。这些案例共同表明两轮迭代机制能够根据第一轮的实际结果自适应地调整分析方向，而非简单重复或放大初始假设。

\subsection{本章小结}

本章通过三个完整案例（技术影响力驱动因素分析、技术融合趋势识别、国家竞争态势评估）展示了系统从“研究问题输入”到“洞察生成—任务规划—代码执行—结果解释”的端到端应用流程。三个案例覆盖因果分析、趋势识别和比较评估三类典型专利计量任务，共产出 15 个分析任务且全部执行成功，说明在当前实现下，数据感知机制和两轮迭代机制已经能够在不同场景中稳定地生成可执行、可追溯的分析工作流。案例同时表明，系统并非简单替代研究者的领域判断，而是以“洞察—任务—证据链”的形式提供结构化决策支持，降低了手工试错成本、提高了方案与数据之间的一致性。需要强调的是，本章侧重展示执行链路和输出形态，本身不构成对系统效果的独立评估；关于各模块的相对贡献及总体优劣，仍应结合第五章的对比实验结果以及后续可能的重复实验一并加以判断。
